%% ma: L10-B14-0010
%% lop: 10
%% chuong: 5
%% bai: 14
%% dang: 10.14.2
%% loai: TL
%% muc: VD
%% dap_an: "-"
%% nguon: "chuyen-de-cac-so-dac-trung-cua-mau-so-lieu-khong-ghep-nhom-toan-10-chuong-trinh-moi.pdf (GV. Phan Nhật Linh) (P04), Bài 03 Dạng 1 Bài tập tự luận 9 (thứ hai)"
%% kien_thuc: "Lớp 10 Bài 13 (tứ phân vị), Bài 14 (khoảng tứ phân vị)"
%% trang_thai: chinh-thuc
%% ghi_chu: "Biểu đồ đường vẽ lại bằng TikZ theo số liệu đọc từ hình gốc (có đường gióng tới từng điểm). Tự duyệt theo yêu cầu giáo viên 22/10/2026."
\begin{ex}
Số lượng muối tiêu thụ (đơn vị: tấn) của một cửa hàng qua các tháng được biểu thị qua biểu đồ sau:
\begin{center}
\begin{tikzpicture}[x=.7cm,y=.0011cm]
  \foreach \y in {1000,2000,3000,4000,5000}{\draw[gray!40] (0,\y)--(12.5,\y); \node[left,font=\scriptsize] at (0,\y) {$\y$};}
  \foreach \m/\v in {1/1500,2/1500,3/2000,4/2000,5/3000,6/3500,7/4000,8/5000,9/4000,10/3000,11/1500,12/1000}{
    \draw[phu] (\m,0)--(\m,\v);
    \node[below,font=\scriptsize] at (\m,0) {$\m$};
  }
  \draw[phu] (0,3500)--(6,3500); \node[left,font=\scriptsize] at (0,3500) {$3500$};
  \draw[thick] (1,1500)--(2,1500)--(3,2000)--(4,2000)--(5,3000)--(6,3500)--(7,4000)--(8,5000)--(9,4000)--(10,3000)--(11,1500)--(12,1000);
  \foreach \m/\v in {1/1500,2/1500,3/2000,4/2000,5/3000,6/3500,7/4000,8/5000,9/4000,10/3000,11/1500,12/1000} \fill (\m,\v) circle (1.6pt);
  \draw[truc] (0,0)--(13,0) node[below]{\footnotesize Tháng};
  \draw[truc] (0,0)--(0,5600) node[above]{\footnotesize Tấn};
  \node[below left,font=\scriptsize] at (0,0) {$O$};
\end{tikzpicture}
\end{center}
Tính khoảng biến thiên và khoảng tứ phân vị của mẫu số liệu.
\loigiai{Đọc biểu đồ: $1500$; $1500$; $2000$; $2000$; $3000$; $3500$; $4000$; $5000$; $4000$; $3000$; $1500$; $1000$.

$R=5000-1000=4000$.

Sắp xếp: $1000$; $1500$; $1500$; $1500$; $2000$; $2000$; $3000$; $3000$; $3500$; $4000$; $4000$; $5000$.

$Q_1=\dfrac{1500+1500}2=1500$; $Q_2=\dfrac{2000+3000}2=2500$; $Q_3=\dfrac{3500+4000}2=3750$.

$\Delta_Q=3750-1500=2250$.}
\end{ex}
